\lecture{2}{Tue 08 Sep 2026 11:02}{Affine varieties}

We now have a course page \href{https://math.bu.edu/people/szczesny/Teaching/745F26/745.html}{here}. We will now be following Galthmann's latest notes \cite{gathmann2021}. Note that these notes define an affine variety as what we previously called an algebraic set.

\begin{definition}[Affine variety]\label{dfn:1-1-6}
	An \emph{affine variety} is an algebraic set (see \cref{dfn:1.1.2}). Some authors (and our original notes) require that affine varieties be \emph{irreducible} algebraic sets, but we do not make that distinction. Thus, we refer only to affine varieties in the sequel.
\end{definition}


Let $k$ be a field, usually algebraically closed. We will switch to denoting the vanishing locus of $I\subseteq k[x_1, \ldots ,x_n]$ by $V(I)$ rather than $Z(I)$. Recall that we topologize affine space $\mathbb{A} _n^k = k^n$ in a nonstandard way, by taking the closed sets to be precisely $Z(I)$ for all ideals $I$ (equivalently subsets $S\subseteq k[x_1, \ldots ,x_n]$, since $V(S)=V((S))$, for $(S)$ the ideal generated by $S$).

\begin{exercise}\label{ex:1-1-7}
	Consider $\mathbb{A} ^3_k$. Let $I$ be the ideal $(x^2+y^2+z^2-1)$. Let $J$ be the ideal $(x,y)$. What is $V(I)\cap V(J)$?
\end{exercise}
\begin{proof}
	We expect $V(I)$ to be the sphere, and we expect $J$ to be the $z$-axis. Thus the intersection should be the two intersection points $(0,0,\pm1)$. By \cref{lma:1.1.4}, we have $V(I)\cap V(J)=V(I+J)$. We then have
	\[
		V(I+J)=V(x,y,x^2+y^2+z^2-1)
	.\]
	Thus we need a point $(x,y,z) = (0,0,z)$ with $z^2-1=0$. This indeed agrees with our geometric intuition (since the field is algebraically closed, yielding 2 solutions).
\end{proof}

Recall that if $R$ is a commutative ring and $I\subseteq R$ is an ideal, then the radical $\sqrt{I} $ is the ideal of all $n$th roots of elements of $I$, defined more precisely as
\[
	\sqrt{I} =\left\{ x\in R \mid x^n\in I \right\} 
.\]
\begin{definition}[Radical]\label{dfn:radical-ideal}
	Let $R$ be a commutative ring, and let $I\subseteq R$ be an ideal. The \emph{radical ideal} of $I$, denoted $\sqrt{I} $, consists of all $n$th roots of elements of $I$. More precisely,
	\[
		\sqrt{I} \coloneqq \left\{ x\in R\mid \exists n\in\mathbb{N} \text{ s.t. } x^n\in I  \right\} 
	.\]
\end{definition}

\begin{exercise}\label{ex:1-1-9}
	Show that
	\[
		\sqrt{I} =\bigcap_{\substack{\mathfrak{p} \text{ prime} \\ I \subseteq \mathfrak{p} }} \mathfrak{p} 
	.\]
\end{exercise}

If $I=0$, then $\sqrt{I}=\sqrt{0}  $ is denotd $\mathfrak{R} $, and called the \emph{nilradical}. It consists of all nilpotent elements of $R$, and by \cref{ex:1-1-8} it follows that $\mathfrak{R} $ is the intersection of all prime ideals of $R$.

Fun fact: if $\pi : R \to R/I$ is the quotient map, then $\sqrt{I} =\pi ^{-1} (\mathfrak{R} )$. Another one: $\sqrt{(x^n)} =(x)$. We also have $\sqrt{(x^2,y^3)} =(x,y)$.

Let $j$ be the map taking ideals in $k[x_1, \ldots ,x_n]$ to closed subsetts in $\mathbb{A} _k^n$ by mapping $I\mapsto V(I)$. We claim that $j(I)=j(\sqrt{I} )$. For example, $V(x^4) = 0 = V(x)$. We can consider the inverse direction. Let $X\subseteq A_n^k$; we can look at the ideal $I(X)$ of all functions $f\in k[x_1, \ldots ,x_n]$ that vanish on $X$. This yields maps
\[
	j : \left\{ \text{ideals of } k[x_1, \ldots ,x_n] \right\} \rightleftarrows \left\{ \text{closed subsets of } \mathbb{A} _k^n \right\} : I
.\]
How do these maps relate? In other words, what are $I(j(J)) = I(V(J))$ and $V(I(X))$ for $J$ an ideal and $X$ a closed subset? This follows by Hilbert's nullstellensatz.

\begin{theorem}[Hilbert's nullstellensatz]\label{thm:1-1-10}
	Let $X\subseteq \mathbb{A} _k^n$ be a closed subset, and let $k$ be an algebraically closed field.
	\begin{enumerate}
		\item $V(\sqrt{I(X)} )=V(I(X))=X$.
		\item $I(X) = \sqrt{I(X)} $.
	\end{enumerate}
\end{theorem}

Consider an example. Let $J\subseteq k[x]$ to be the ideal $J=((x-a_1)^{n_1} \cdots (x-a_r)^{n_r} )$, for $n_j>0$ and $a_i \neq a_j$ when $i\neq j$. Write $f$ for this polynomial, such that $J=(f)$. Since $k[x]$ is a principal ideal, $\sqrt{J} $ is generated by a single element. We have
\[
	\sqrt{J} =((x-a_1)\cdots (x-a_r))
.\]
We see that the vanishing loci of $J$ and $\sqrt{J} $ are clearly equal.


\begin{remark}\label{rmk:1-1-11}
	Recall by \cref{ex:1-1-8} that all radicals are intersections of primes. Thus, any prime ideal $\mathfrak{p} $ is radical. Moreover, all maximal ideals are radical. Given affine varieties $X_1,X_2 \subseteq \mathbb{A} _k^n$, we can even show
	\[
		I(X_1 \cup X_2) = I(X_1)\cap I(X_2),\qquad I(X_1 \cap X_2) = \sqrt{I(X_1)+I(X_2)} 
	.\]
	(exercise: how to prove this?) This implies that interections of radicals are radical.
\end{remark}
